% GENERATED FILE. Edit canonical lessons, then run npm run generate:books.
% canonical-source-hash: fnv1a64:1636c7c34405bde5
% canonical-lessons: HI-C402-gathri, HI-C402-ataici, HI-C402-sarkana, HI-C402-ughna, HI-C402-sustana

\chapter{Bundles and Naps}
\label{ch:hi-bundles-and-naps}
\hlchaptermodality{402}

\begin{chapteropening}
\textbf{By the end of this chapter:} \emph{I can stow my bags, carry a bundle and doze on the way.}

Complete the last lesson of chapter 402: I can stow my bags, carry a bundle and doze on the way.
\end{chapteropening}

\section[\texorpdfstring{\hi{गठरी} (gaṭhrī)}{गठरी (gaṭhrī)}]{\hi{गठरी} (gaṭhrī) — a bundle (tied in cloth)}

\label{lesson:HI-C402-gathri}



\begin{quote}
Before the new one: say the Hindi for to step over, then the Hindi for to look out (at).
\end{quote}

\subsection*{You'll want to know: \hi{गठरी}}
\textbf{\hi{गठरी}} — \emph{gaṭhrī} — \textquotedblleft{}a bundle (tied in cloth)\textquotedblright{}.

\begin{quote}
\hi{दादी} \hi{के} \hi{पास} \hi{कपड़ों} \hi{की} \hi{एक} \hi{गठरी} \hi{थी।}
\end{quote}

\begin{quote}
\emph{dādī ke pās kapṛoṁ kī ek gaṭhrī thī.}
\end{quote}

\begin{quote}
Grandmother had a bundle of clothes.
\end{quote}

Feminine.

\begin{cousinweb}
\textbf{\hi{गठरी}} (\emph{gaṭhrī}) is a tadbhava of Sanskrit \emph{grathna}, \textquotedblleft{}a bunch\textquotedblright{} (CDIAL 4348), from the root of tying knots; CDIAL lists Hindi \emph{gaṭhṛī}, \textquotedblleft{}bundle\textquotedblright{}.
\end{cousinweb}

\begin{grammarlens}[title={What you've built}]
One more word for getting around.
\end{grammarlens}

\subsection*{Your turn}
\begin{itemize}
\raggedright
  \item \emph{Say it:} \emph{gaṭhrī}
  \item \emph{Say it:} \emph{gaṭhrī}, once more
  \item \emph{Say it:} \emph{dādī ke pās kapṛoṁ kī ek gaṭhrī thī.} — Grandmother had a bundle of clothes.
  \item \emph{Recall:} say the Hindi for to step over, then the Hindi for to look out (at), then say \emph{gaṭhrī} again
\end{itemize}

\begin{tcolorbox}[breakable,colback=teal!4,colframe=teal!35!black,title={Before you move on}]
What does \hi{गठरी} mean? (\textquotedblleft{}A bundle (tied in cloth)\textquotedblright{}.) Say it once more.
\end{tcolorbox}
\section[\texorpdfstring{\hi{अटैची} (aṭaicī)}{अटैची (aṭaicī)}]{\hi{अटैची} (aṭaicī) — a suitcase}

\label{lesson:HI-C402-ataici}



\begin{quote}
Before the new one: say the Hindi for to look out (at), then the Hindi for a bundle.
\end{quote}

\subsection*{You'll want to know: \hi{अटैची}}
\textbf{\hi{अटैची}} — \emph{aṭaicī} — \textquotedblleft{}a suitcase\textquotedblright{}.

\begin{quote}
\hi{मेरी} \hi{अटैची} \hi{बहुत} \hi{भारी} \hi{है।}
\end{quote}

\begin{quote}
\emph{merī aṭaicī bahut bhārī hai.}
\end{quote}

\begin{quote}
My suitcase is very heavy.
\end{quote}

Feminine.

\begin{cousinweb}
\textbf{\hi{अटैची}} (\emph{aṭaicī}) is borrowed from English \emph{attaché}, short for \emph{attaché case}; in India it names any small suitcase, and it is feminine.
\end{cousinweb}

\begin{grammarlens}[title={What you've built}]
One more word for getting around.
\end{grammarlens}

\subsection*{Your turn}
\begin{itemize}
\raggedright
  \item \emph{Say it:} \emph{aṭaicī}
  \item \emph{Say it:} \emph{aṭaicī}, once more
  \item \emph{Say it:} \emph{merī aṭaicī bahut bhārī hai.} — My suitcase is very heavy.
  \item \emph{Recall:} say the Hindi for to look out (at), then the Hindi for a bundle, then say \emph{aṭaicī} again
\end{itemize}

\begin{tcolorbox}[breakable,colback=teal!4,colframe=teal!35!black,title={Before you move on}]
What does \hi{अटैची} mean? (\textquotedblleft{}A suitcase\textquotedblright{}.) Say it once more.
\end{tcolorbox}
\section[\texorpdfstring{\hi{सरकाना} (sarkānā)}{सरकाना (sarkānā)}]{\hi{सरकाना} (sarkānā) — to slide (something) along; to push aside}

\label{lesson:HI-C402-sarkana}



\begin{quote}
Before the new one: say the Hindi for a bundle, then the Hindi for a suitcase.
\end{quote}

\subsection*{You'll want to know: \hi{सरकाना}}
\textbf{\hi{सरकाना}} — \emph{sarkānā} — \textquotedblleft{}to slide (something) along; to push aside\textquotedblright{}.

\begin{quote}
\hi{मैंने} \hi{अपनी} \hi{अटैची} \hi{सीट} \hi{के} \hi{नीचे} \hi{सरकाई।}
\end{quote}

\begin{quote}
\emph{maiṁne apnī aṭaicī sīṭ ke nīce sarkāī.}
\end{quote}

\begin{quote}
I slid my suitcase under the seat.
\end{quote}

Past, with \hi{ने} (he or she …): \hi{उसने} \hi{कुर्सी} \hi{पीछे} \hi{सरकाई}, \emph{usne kursī pīche sarkāī}.

\begin{cousinweb}
\textbf{\hi{सरकाना}} (\emph{sarkānā}) is the causative partner of \hi{सरकना}, from Sanskrit \emph{sarati}, \textquotedblleft{}glides\textquotedblright{} (CDIAL 13250): \hi{सरकना} is to slide, \hi{सरकाना} to make a thing slide.
\end{cousinweb}

\begin{grammarlens}[title={What you've built}]
One more word for getting around.
\end{grammarlens}

\subsection*{Your turn}
\begin{itemize}
\raggedright
  \item \emph{Say it:} \emph{sarkānā}
  \item \emph{Say it:} \emph{sarkānā}, once more
  \item \emph{Say it:} \emph{maiṁne apnī aṭaicī sīṭ ke nīce sarkāī.} — I slid my suitcase under the seat.
  \item \emph{Recall:} say the Hindi for a bundle, then the Hindi for a suitcase, then say \emph{sarkānā} again
\end{itemize}

\begin{tcolorbox}[breakable,colback=teal!4,colframe=teal!35!black,title={Before you move on}]
What does \hi{सरकाना} mean? (\textquotedblleft{}To slide (something) along; to push aside\textquotedblright{}.) Say it once more.
\end{tcolorbox}
\section[\texorpdfstring{\hi{ऊँघना} (ū̃ghnā)}{ऊँघना (ū̃ghnā)}]{\hi{ऊँघना} (ū̃ghnā) — to doze}

\label{lesson:HI-C402-ughna}



\begin{quote}
Before the new one: say the Hindi for a suitcase, then the Hindi for to slide (something).
\end{quote}

\subsection*{You'll want to know: \hi{ऊँघना}}
\textbf{\hi{ऊँघना}} — \emph{ū̃ghnā} — \textquotedblleft{}to doze\textquotedblright{}.

\begin{quote}
\hi{बस} \hi{में} \hi{बहुत} \hi{लोग} \hi{ऊँघ} \hi{रहे} \hi{थे।}
\end{quote}

\begin{quote}
\emph{bas meṁ bahut log ū̃gh rahe the.}
\end{quote}

\begin{quote}
A lot of people on the bus were dozing.
\end{quote}

Past (he …): \hi{वह} \hi{थोड़ी} \hi{देर} \hi{ऊँघा}, \emph{vah thoṛī der ū̃ghā}.

\begin{cousinweb}
\textbf{\hi{ऊँघना}} (\emph{ū̃ghnā}) goes back to \emph{*uṅghati}, \textquotedblleft{}sleeps\textquotedblright{}, a form CDIAL reconstructs (CDIAL 1632), which lists Hindi \emph{ū̃ghnā}, \textquotedblleft{}to be drowsy\textquotedblright{}.
\end{cousinweb}

\begin{grammarlens}[title={What you've built}]
One more word for getting around.
\end{grammarlens}

\subsection*{Your turn}
\begin{itemize}
\raggedright
  \item \emph{Say it:} \emph{ū̃ghnā}
  \item \emph{Say it:} \emph{ū̃ghnā}, once more
  \item \emph{Say it:} \emph{bas meṁ bahut log ū̃gh rahe the.} — A lot of people on the bus were dozing.
  \item \emph{Recall:} say the Hindi for a suitcase, then the Hindi for to slide (something), then say \emph{ū̃ghnā} again
\end{itemize}

\begin{tcolorbox}[breakable,colback=teal!4,colframe=teal!35!black,title={Before you move on}]
What does \hi{ऊँघना} mean? (\textquotedblleft{}To doze\textquotedblright{}.) Say it once more.
\end{tcolorbox}
\section[\texorpdfstring{\hi{सुस्ताना} (sustānā)}{सुस्ताना (sustānā)}]{\hi{सुस्ताना} (sustānā) — to rest a while}

\label{lesson:HI-C402-sustana}



\begin{quote}
Before the new one: say the Hindi for to slide (something), then the Hindi for to doze.
\end{quote}

\subsection*{You'll want to know: \hi{सुस्ताना}}
\textbf{\hi{सुस्ताना}} — \emph{sustānā} — \textquotedblleft{}to rest a while\textquotedblright{}.

\begin{quote}
\hi{हम} \hi{पेड़} \hi{के} \hi{नीचे} \hi{थोड़ी} \hi{देर} \hi{सुस्ताए।}
\end{quote}

\begin{quote}
\emph{ham peṛ ke nīce thoṛī der sustāe.}
\end{quote}

\begin{quote}
We rested for a while under the tree.
\end{quote}

Past (he …): \hi{वह} \hi{पेड़} \hi{के} \hi{नीचे} \hi{सुस्ताया}, \emph{vah peṛ ke nīce sustāyā}.

\begin{cousinweb}
\textbf{\hi{सुस्ताना}} (\emph{sustānā}) is built on \emph{sust}, \textquotedblleft{}slack, weary\textquotedblright{}, usually traced to Persian \emph{sust}, with the Hindi verb ending -ānā: to let your tiredness go for a while.
\end{cousinweb}

\begin{grammarlens}[title={What you've built}]
One more word for getting around.
\end{grammarlens}

\subsection*{Your turn}
\begin{itemize}
\raggedright
  \item \emph{Say it:} \emph{sustānā}
  \item \emph{Say it:} \emph{sustānā}, once more
  \item \emph{Say it:} \emph{ham peṛ ke nīce thoṛī der sustāe.} — We rested for a while under the tree.
  \item \emph{Recall:} say the Hindi for to slide (something), then the Hindi for to doze, then say \emph{sustānā} again
\end{itemize}

\begin{tcolorbox}[breakable,colback=teal!4,colframe=teal!35!black,title={Before you move on}]
What does \hi{सुस्ताना} mean? (\textquotedblleft{}To rest a while\textquotedblright{}.) Say it once more.
\end{tcolorbox}
